\chapter{Geluid en akoestiek}\label{ch:g12:sound-acoustics}

Tokkel de lage snaar van een gitaar: de kamer vult zich met een lage a, en
toch is er niets naar je oor gereisd behalve samengeperste en uitgerekte
lucht. Dit hoofdstuk volgt dat patroon --- zijn snelheid, hoe \omterm{def:g10:signals-and-waves:frequency}{frequentie}
\omterm{def:g12:sound-acoustics:pitch-loudness}{toonhoogte} wordt en \omterm{def:g11:work-of-force:power}{vermogen} \omterm{def:g12:sound-acoustics:pitch-loudness}{luidheid}, en waarom één noot op elk
instrument anders klinkt.

\section{Geluid is een drukgolf}

\begin{definition}[Geluidsgolf]\label{def:g12:sound-acoustics:sound-wave}
Een \emph{geluidsgolf}\index{geluidsgolf} is een mechanische drukgolf
(\cref{ch:g12:mechanical-waves}): een trillend oppervlak wekt reizende
\emph{verdichtingen}\index{verdichting} (lichte \omterm{thm:g10:pressure:depth}{overdruk}) en
\emph{verdunningen}\index{verdunning} (onderdruk) op. Luchtpakketjes
trillen \emph{langs} de looprichting (\omterm{def:g10:signals-and-waves:sound}{geluid} is
\emph{longitudinaal}\index{longitudinale golf}), en er is een tussenstof
nodig: een bel onder een vacuümstolp verstomt
(\cref{ch:g10:signals-and-waves}).
\end{definition}

\begin{omfigure}
\begin{tikzpicture}[font=\small]
  \fill[gray!55] (-0.75,-0.4) rectangle (-0.35,0.4);
  \draw[gray!75, fill=gray!20] (-0.35,0.28) -- (0.0,0.5) -- (0.0,-0.5)
    -- (-0.35,-0.28) -- cycle;
  \foreach \i in {0,...,69} \draw[gray!80, thin]
    ({0.2+\i*0.1+0.12*sin(\i*12.857)},-0.45) -- +(0,0.9);
  \draw[dashed, thin] (1.6,0.5) -- (1.6,0.7) (4.4,0.5) -- (4.4,0.7);
  \draw[omThm, Stealth-Stealth] (1.6,0.7) -- (4.4,0.7)
    node[midway, above] {$\lambda$};
  \draw[-Stealth] (1.6,-0.55) -- (1.6,-0.95) node[below] {verdichting};
  \draw[-Stealth] (3.0,-0.55) -- (3.0,-0.6) node[below] {verdunning};
  \draw[-Stealth, omDef, thick] (7.45,0) -- (8.15,0) node[right] {$v$};
\end{tikzpicture}

{\small Een luidspreker drijft een \omterm{def:g12:sound-acoustics:sound-wave}{longitudinale golf} aan: pakketjes
hopen zich op tot \omterm{def:g12:sound-acoustics:sound-wave}{verdichtingen} die één \omterm{def:g12:mechanical-waves:wavelength}{golflengte}
$\lambda = v/f$ uit elkaar liggen.}
\end{omfigure}

\begin{proposition}[Geluidssnelheid]\label{prop:g12:sound-acoustics:speed}
\index{geluidssnelheid}%
\omterm{def:g10:signals-and-waves:sound}{Geluid} reist met ongeveer \qty{340}{m/s} in lucht van
\qty{20}{\degreeCelsius}, met \qty{1500}{m/s} in water en met
\qty{5000}{m/s} in staal: hoe stijver het midden, hoe sneller.
\admitted
\end{proposition}

\begin{remark}[Warme lucht is sneller]\label{rem:g12:sound-acoustics:temperature}
Dit zijn gemeten waarden; ze voorspellen gebeurt in het volume van
bachelorjaar 1. \omterm{def:g10:signals-and-waves:sound}{Geluid} gaat sneller in warmere lucht (\qty{331}{m/s} bij
\qty{0}{\degreeCelsius}) --- en daarom raken blaasinstrumenten ontstemd
naarmate de zaal opwarmt.
\end{remark}

\begin{example}[De afmetingen van geluiden]\label{ex:g12:sound-acoustics:wavelengths}
De betrekking $\lambda = v/f$ geeft elk \omterm{def:g10:signals-and-waves:sound}{geluid} zijn maat: in lucht
overspant een gerommel van \qty{20}{Hz} $340/20 = \qty{17}{m}$ --- een
gebouw --- terwijl een gesis van \qty{20}{kHz} in \qty{1.7}{cm} past; de
metergrote buigen makkelijk om deuren en hoeken heen.
\end{example}

\section{Toonhoogte, luidheid en het oor}

\begin{definition}[Toonhoogte en luidheid]\label{def:g12:sound-acoustics:pitch-loudness}
\emph{Toonhoogte}\index{toonhoogte} is de waarneming van de \omterm{def:g10:signals-and-waves:frequency}{frequentie}:
hoe hoger $f$, hoe hoger de noot; de kamertoon a is \qty{440}{Hz}, en $f$
verdubbelen tilt de noot een octaaf hoger.
\emph{Luidheid}\index{luidheid} is de waarneming van de \omterm{def:g10:signals-and-waves:amplitude}{amplitude}: hoe
sterker de drukwisseling, hoe luider; versterken verandert de luidheid,
niet de toonhoogte.
\end{definition}

\begin{remark}[Het venster van het oor]\label{rem:g12:sound-acoustics:audible}
Het gehoor loopt van ongeveer \qty{20}{Hz} tot \qty{20}{kHz}
(\cref{ch:g10:signals-and-waves}); daaronder ligt \omterm{def:g10:signals-and-waves:ultrasound}{infrageluid},
daarboven \omterm{def:g10:signals-and-waves:ultrasound}{ultrageluid} --- vleermuizen op \qty{50}{kHz}, medische sondes op
megahertz. Een piano bestrijkt alleen het midden:
\qtyrange{27.5}{4186}{Hz}.
\end{remark}

\section{Geluidsintensiteit en de decibelschaal}

\begin{definition}[Geluidsintensiteit]\label{def:g12:sound-acoustics:intensity}
De \emph{geluidsintensiteit}\index{geluidsintensiteit} $I$ is het
akoestische \omterm{def:g11:work-of-force:power}{vermogen} dat door één vierkante \omterm{def:g10:orders-of-magnitude:unit}{meter} loodrecht op de \omterm{def:g10:signals-and-waves:wave}{golf}
gaat: $I = P/S$, in \unit{W/m^2}. Het zwakste hoorbare \omterm{def:g10:signals-and-waves:sound}{geluid} is ongeveer
$I_0 = \qty{1.0e-12}{W/m^2}$; de pijn begint rond \qty{1}{W/m^2}.
\end{definition}

\begin{proposition}[Verzwakking met het omgekeerde kwadraat]\label{prop:g12:sound-acoustics:inverse-square}
\index{omgekeerd kwadratische wet}%
Een kleine bron die een \omterm{def:g11:work-of-force:power}{vermogen} $P$ in alle richtingen even sterk
uitstraalt, geeft op afstand $r$ de intensiteit
$I = P/(4\pi r^2) \propto 1/r^2$: de afstand verdubbelen deelt de
intensiteit door vier.
\end{proposition}

\begin{proof}
Het hele \omterm{def:g11:work-of-force:power}{vermogen} $P$ gaat door de bol met oppervlakte $4\pi r^2$; deel.
\end{proof}

\begin{definition}[Geluidsniveau]\label{def:g12:sound-acoustics:level}
Het oor bestrijkt twaalf machten van tien, dus worden intensiteiten
logaritmisch opgegeven. Het \emph{geluidsniveau}\index{geluidsniveau} van
$I$ is
\[ L = 10 \log_{10}(I/I_0), \qquad I_0 = \qty{1.0e-12}{W/m^2}, \]
in \emph{decibel}\index{decibel} (\unit{dB}): de gehoordrempel ligt op
\qty{0}{dB}, de pijngrens rond \qty{120}{dB}.
\end{definition}

\begin{proposition}[Rekenen met decibel]\label{prop:g12:sound-acoustics:db-arithmetic}
$I$ met $k$ vermenigvuldigen telt er $10\log_{10}k$ \omterm{def:g12:sound-acoustics:level}{decibel} bij. In het
bijzonder voegt $\times 10$ er \qty{10}{dB} aan toe, en $\times 2$
$10\log_{10}2 \approx \qty{3.0}{dB}$; $N$ onafhankelijke gelijke bronnen
tellen hun \emph{intensiteiten} op, nooit hun niveaus, en geven
$L_1 + 10\log_{10}N$; de afstand tot een kleine bron verdubbelen ($I$
gedeeld door vier) haalt er ongeveer \qty{6.0}{dB} af.
\end{proposition}

\begin{proof}
$10\log_{10}(kI/I_0) = 10\log_{10}(I/I_0) + 10\log_{10}k$; voor de
afstand: \cref{prop:g12:sound-acoustics:inverse-square} met $k = 1/4$.
\end{proof}

\begin{method}[Werken met decibel]\label{met:g12:sound-acoustics:db}
\begin{enumerate}
  \item $L = 10\log_{10}(I/I_0)$; terug: $I = I_0 \times 10^{L/10}$.
  \item Tel nooit niveaus op: zet om naar intensiteiten, tel op, en zet
        terug --- of gebruik $+\qty{3}{dB}$ per verdubbeling en
        $+\qty{10}{dB}$ per vertienvoudiging.
  \item Elke verdubbeling van de afstand tot een kleine bron kost
        \qty{6}{dB}.
\end{enumerate}
\end{method}

\begin{omfigure}
\begin{tikzpicture}[font=\small]
  \fill[omThm!12] (4.25,-0.14) rectangle (7.35,0.14);
  \draw[-Stealth] (-0.3,0) -- (7.7,0);
  \foreach \x/\l in {0/0, 1/20, 2/40, 3/60, 4/80, 5/100, 6/120, 7/140}
    \draw (\x,0.07) -- (\x,-0.07) node[below] {\l};
  \node[below] at (7.55,-0.07) {\unit{dB}};
  \foreach \x/\y/\t in {0/0.4/gehoordrempel, 1.5/0.95/gefluister,
      3/0.4/gesprek, 4/0.95/{drukke straat}, 5.25/0.4/concert,
      6/0.95/pijngrens, 7/0.4/{opstijgend vliegtuig}}
    \draw[omDef] (\x,0.14) -- (\x,\y) node[above] {\t};
  \draw[omThm, thick] (4.25,-0.14) -- (4.25,0.14);
\end{tikzpicture}

{\small Een ladder van alledaagse niveaus: elke stap van \qty{20}{dB} is
een honderdvoudige sprong in intensiteit; voorbij \qty{85}{dB} (rode zone)
beschadigt blootstelling het gehoor.}
\end{omfigure}

\begin{remark}[Gehoorschade]\label{rem:g12:sound-acoustics:damage}
Gehoorverlies is een kwestie van \omterm{rem:g11:nucleus-radioactivity:dose}{dosis}: \qty{85}{dB} is acht uur lang
veilig, en elke \qty{3}{dB} meer --- een verdubbeling van de intensiteit
--- halveert de veilige tijd; verloren haarcellen groeien niet terug.
\end{remark}

\section{Klankkleur en boventonen}

\begin{definition}[Boventonen]\label{def:g12:sound-acoustics:harmonics}
Een \omterm{def:g10:signals-and-waves:period}{periodiek} \omterm{def:g10:signals-and-waves:sound}{geluid} met \omterm{def:g10:signals-and-waves:frequency}{frequentie} $f_1$ --- de
\emph{grondtoon}\index{grondtoon} --- is in het algemeen een som van
sinussen met de \omterm{def:g10:signals-and-waves:frequency}{frequenties} $f_n = n f_1$, $n = 1, 2, 3, \dots$, zijn
\emph{boventonen}\index{boventoon}. Een \omterm{def:g10:signals-and-waves:sound}{geluid} dat tot zijn grondtoon is
teruggebracht, is een \emph{zuivere toon}\index{zuivere toon}, zoals die
van een stemvork.
\end{definition}

\begin{definition}[Spectrum en klankkleur]\label{def:g12:sound-acoustics:timbre}
Het \emph{spectrum}\index{spectrum} van een \omterm{def:g10:signals-and-waves:sound}{geluid} is het staafdiagram van
zijn \omterm{def:g12:sound-acoustics:harmonics}{boventonen}: één staaf per $f_n$, waarvan de hoogte de sterkte van die
\omterm{def:g12:sound-acoustics:harmonics}{boventoon} is. Het oor hoort $f_1$ als de \omterm{def:g12:sound-acoustics:pitch-loudness}{toonhoogte} en het \emph{mengsel}
van sterkten als de \emph{klankkleur}\index{klankkleur} --- de kleur van
het \omterm{def:g10:signals-and-waves:sound}{geluid}: instrumenten die dezelfde noot spelen delen $f_1$ en
verschillen alleen in hun spectra.
\end{definition}

\begin{omfigure}
\begin{tikzpicture}[font=\small]
  \node at (1.45,1.0) {fluit}; \node at (5.85,1.0) {viool};
  \draw[-Stealth] (-0.1,0) -- (3.1,0) node[below] {$t$};
  \draw[omDef, thick, smooth, samples=120, domain=0:2.9]
    plot (\x, {0.55*sin(248*\x)});
  \begin{scope}[xshift=4.4cm]
    \draw[-Stealth] (-0.1,0) -- (3.1,0) node[below] {$t$};
    \draw[omDef, thick, smooth, samples=160, domain=0:2.9]
      plot (\x, {0.30*sin(248*\x)+0.18*sin(496*\x)
        +0.12*sin(744*\x)+0.08*sin(992*\x)});
  \end{scope}
  \begin{scope}[yshift=-3.3cm]
    \draw[-Stealth] (-0.1,0) -- (3.1,0) node[below] {$f$ (\unit{Hz})};
    \foreach \k/\h in {1/1.15, 2/0.10, 3/0.05} \fill[omProp!80]
      ({0.35+0.55*(\k-1)-0.09},0) rectangle ({0.35+0.55*(\k-1)+0.09},\h);
    \foreach \k/\l in {1/220, 2/440, 3/660, 4/880, 5/1100}
      \node[below] at ({0.35+0.55*(\k-1)},0) {\scriptsize \l};
    \begin{scope}[xshift=4.4cm]
      \draw[-Stealth] (-0.1,0) -- (3.1,0) node[below] {$f$ (\unit{Hz})};
      \foreach \k/\h in {1/0.80, 2/0.55, 3/0.38, 4/0.26, 5/0.18}
        \fill[omProp!80] ({0.35+0.55*(\k-1)-0.09},0)
          rectangle ({0.35+0.55*(\k-1)+0.09},\h);
      \foreach \k/\l in {1/220, 2/440, 3/660, 4/880, 5/1100}
        \node[below] at ({0.35+0.55*(\k-1)},0) {\scriptsize \l};
    \end{scope}
  \end{scope}
\end{tikzpicture}

{\small Dezelfde a (\qty{220}{Hz}): golfvormen (boven), spectra (onder)
--- dezelfde \omterm{def:g12:sound-acoustics:harmonics}{grondtoon} en dezelfde \omterm{def:g12:sound-acoustics:pitch-loudness}{toonhoogte}; de boventonenladder van de
viool is haar \omterm{def:g12:sound-acoustics:timbre}{klankkleur}.}
\end{omfigure}

\begin{example}[Waarom je de beller herkent]\label{ex:g12:sound-acoustics:two-instruments}
De a van een fluit is bijna een \omterm{def:g12:sound-acoustics:harmonics}{zuivere toon}; een viool die dezelfde
\qty{220}{Hz} strijkt voegt er sterke \omterm{def:g12:sound-acoustics:harmonics}{boventonen} op
\qtylist{440;660;880}{Hz} aan toe. Dezelfde \omterm{def:g12:sound-acoustics:pitch-loudness}{toonhoogte} en dezelfde
\omterm{def:g12:sound-acoustics:pitch-loudness}{luidheid}, en toch verwart niemand de twee: ook een stem is een \omterm{def:g12:sound-acoustics:timbre}{spectrum}
dat je uit je hoofd kent.
\end{example}

\section{De trillende snaar}

\begin{definition}[Staande golf]\label{def:g12:sound-acoustics:standing-wave}
Een \omterm{def:g10:signals-and-waves:wave}{golf} die tussen de vaste uiteinden van een snaar opgesloten zit, kan
zich vestigen als \emph{staande golf}\index{staande golf}: een trilling
ter plaatse, waarin roerloze punten --- \emph{knopen}\index{knoop} ---
afwisselen met punten van maximale uitslag --- \emph{buiken}\index{buik}.
Er reist niets meer; de snaar trilt in een vast patroon.
\end{definition}

\begin{proposition}[Trillingswijzen van een snaar]\label{prop:g12:sound-acoustics:string-modes}
Een snaar met lengte $L$, aan beide uiteinden vastgezet, waarin golven met
snelheid $v$ lopen, trilt alleen gestaag in patronen waarin een geheel
aantal $n$ halve \omterm{def:g12:mechanical-waves:wavelength}{golflengten} tussen de opgelegde \omterm{def:g12:sound-acoustics:standing-wave}{knopen} aan de uiteinden
past:
\[ \lambda_n = \frac{2L}{n}, \quad f_n = n\,\frac{v}{2L} = n f_1, \quad n = 1, 2, 3, \dots \]
De laagste trillingswijze $f_1 = v/(2L)$ is de \omterm{def:g12:sound-acoustics:harmonics}{grondtoon}; de andere zijn
precies haar \omterm{def:g12:sound-acoustics:harmonics}{boventonen}. \admitted
\end{proposition}

\begin{remark}[Waar het eerlijke bewijs staat]\label{rem:g12:sound-acoustics:honest}
Dat opgesloten golven precies deze trillingswijzen kiezen, volgt uit het
superponeren van de \omterm{def:g10:signals-and-waves:wave}{golf} met haar weerkaatsingen, een berekening die in
het volume van bachelorjaar 1 wordt uitgevoerd. De meetkunde overtuigt:
beide uiteinden moeten stil blijven, dus moet er een geheel aantal halve
\omterm{def:g12:mechanical-waves:wavelength}{golflengten} in passen.
\end{remark}

\begin{omfigure}
\begin{tikzpicture}[font=\small]
  \foreach \n/\y/\l in {1/0/{$f_1 = \dfrac{v}{2L}$},
      2/-1.6/{$f_2 = 2f_1$}, 3/-3.2/{$f_3 = 3f_1$}} {
    \draw[dashed, omExo] (0,\y) -- (4.6,\y);
    \draw[omDef, thick, smooth, samples=80, domain=0:4.6]
      plot (\x, {\y+0.5*sin(\n*39.13*\x)});
    \draw[omDef!50, thick, smooth, samples=80, domain=0:4.6]
      plot (\x, {\y-0.5*sin(\n*39.13*\x)});
    \fill (0,\y) circle (2pt) (4.6,\y) circle (2pt);
    \node[left] at (-0.25,\y) {$n = \n$};
    \node[right] at (4.95,\y) {\l};
  }
  \fill[omThm] (2.3,-1.6) circle (1.8pt);
  \fill[omThm] (1.533,-3.2) circle (1.8pt) (3.067,-3.2) circle (1.8pt);
  \node[omThm, below] at (2.3,-1.72) {knoop};
\end{tikzpicture}

{\small De eerste drie trillingswijzen van een aan beide kanten
vastgezette snaar: er passen $n$ halve \omterm{def:g12:mechanical-waves:wavelength}{golflengten} in $L$, dus
$\lambda_n = 2L/n$ en $f_n = nf_1$.}
\end{omfigure}

\begin{example}[De a-snaar van een viool]\label{ex:g12:sound-acoustics:violin}
De a-snaar van een viool heeft $L = \qty{32.8}{cm}$ en klinkt op
$f_1 = \qty{440}{Hz}$: golven lopen er met
$v = 2Lf_1 = 2 \times 0.328 \times 440 \approx \qty{289}{m/s}$ overheen
--- bepaald door \omterm{def:g11:circuits-and-power:voltage}{spanning} en massa, gestemd met de stemschroef. Een vinger
erop drukken verkort $L$ en tilt elke $f_n$ tegelijk op: het \omterm{def:g12:sound-acoustics:timbre}{spectrum}
schuift omhoog.
\end{example}

\begin{remark}[Waarom instrumenten verschillen]\label{rem:g12:sound-acoustics:instruments}
Hoe de snaar wordt aangeslagen --- getokkeld, gestreken, aangeslagen ---
bepaalt hoeveel er van elke trillingswijze aanwezig is, en de klankkast
versterkt sommige \omterm{def:g12:sound-acoustics:harmonics}{boventonen} meer dan andere; aanslag plus klankkast
leggen het \omterm{def:g12:sound-acoustics:timbre}{spectrum} vast, zodat dezelfde $f_1$ een gitaar, een viool en
een piano elk als zichzelf laat klinken.
\end{remark}

\section{Oefeningen}

\begin{exercise}[$\star$]\label{exo:g12:sound-acoustics:1}
Bepaal in lucht ($v = \qty{340}{m/s}$) de \omterm{def:g12:mechanical-waves:wavelength}{golflengte} van (a) een basnoot
van \qty{55}{Hz}; (b) een piepje van \qty{1.0}{kHz}; (c) een fluittoon van
\qty{15}{kHz}. Vergelijk de grootste met de kleinste.
\end{exercise}

\begin{exercise}[$\star$]\label{exo:g12:sound-acoustics:2}
Deel in bij \omterm{def:g10:signals-and-waves:ultrasound}{infrageluid}, hoorbaar \omterm{def:g10:signals-and-waves:sound}{geluid} of \omterm{def:g10:signals-and-waves:ultrasound}{ultrageluid}: \qty{8}{Hz};
\qty{100}{Hz}; \qty{15}{kHz}; \qty{40}{kHz}; \qty{3}{MHz}. Welke van de
hoorbare klinkt lager van \omterm{def:g12:sound-acoustics:pitch-loudness}{toonhoogte}?
\end{exercise}

\begin{exercise}[$\star$]\label{exo:g12:sound-acoustics:3}
Een werker slaat op een rail \qty{1.0}{km} verderop. Bereken de
looptijden door het staal en door de lucht. Wat hoort iemand die één oor
tegen de rail houdt?
\end{exercise}

\begin{exercise}[$\star$]\label{exo:g12:sound-acoustics:4}
(a) Wat is het \omterm{def:g12:sound-acoustics:level}{geluidsniveau} bij $I = \qty{1.0e-5}{W/m^2}$? (b) Welke
intensiteit hoort bij \qty{60}{dB}? (c) Noem bij elk niveau een
alledaags \omterm{def:g10:signals-and-waves:sound}{geluid}.
\end{exercise}

\begin{exercise}[$\star$]\label{exo:g12:sound-acoustics:5}
Een g-snaar van een viool klinkt op $f_1 = \qty{196}{Hz}$: geef haar
volgende drie \omterm{def:g12:sound-acoustics:harmonics}{boventonen}. Waarom bereikt een fluittoon van \qty{12}{kHz},
hoe rijk ook bij de bron, het oor als een \omterm{def:g12:sound-acoustics:harmonics}{zuivere toon}?
\end{exercise}

\begin{exercise}[$\star\star$]\label{exo:g12:sound-acoustics:6}
Een sirene straalt $P = \qty{0.50}{W}$ in alle richtingen even sterk uit.
Bereken de intensiteit en het \omterm{def:g12:sound-acoustics:level}{geluidsniveau} op \qty{5.0}{m} afstand.
\end{exercise}

\begin{exercise}[$\star\star$]\label{exo:g12:sound-acoustics:7}
Eén viool geeft \qty{68}{dB} op jouw plaats. Welk niveau geven er twee?
En vier? Hoeveel violen zijn er nodig om \qty{78}{dB} te halen?
\end{exercise}

\begin{exercise}[$\star\star$]\label{exo:g12:sound-acoustics:8}
Op \qty{2.0}{m} van een kleine luidspreker is het niveau \qty{95}{dB}.
Voorspel het niveau op \qty{8.0}{m}: onder de grens van \qty{85}{dB}?
\end{exercise}

\begin{exercise}[$\star\star$]\label{exo:g12:sound-acoustics:9}
Een noot van \qty{440}{Hz} gaat van lucht over in water. Bereken haar
\omterm{def:g12:mechanical-waves:wavelength}{golflengte} in beide middens. De duiker hoort dezelfde \omterm{def:g12:sound-acoustics:pitch-loudness}{toonhoogte}: welke
grootheid ligt door de bron vast, en welke past zich aan?
\end{exercise}

\begin{exercise}[$\star\star$]\label{exo:g12:sound-acoustics:10}
Een \omterm{def:g12:sound-acoustics:timbre}{spectrum} toont pieken op \qtylist{130;260;390;520}{Hz} met afnemende
hoogten: wat is de \omterm{def:g12:sound-acoustics:harmonics}{grondtoon}? Een ander instrument speelt dezelfde noot
even luid: wat is er hetzelfde aan zijn \omterm{def:g12:sound-acoustics:timbre}{spectrum}, en wat verschilt?
\end{exercise}

\begin{exercise}[$\star\star$]\label{exo:g12:sound-acoustics:11}
Een gitaarsnaar van \qty{65.0}{cm} klinkt op $f_1 = \qty{110}{Hz}$.
Bepaal de golfsnelheid, en daarna de trillende lengten die het octaaf
(\qty{220}{Hz}) en de kwint (\qty{165}{Hz}) opleveren.
\end{exercise}

\begin{exercise}[$\star\star\star$]\label{exo:g12:sound-acoustics:12}
Een medische sonde stuurt \omterm{def:g10:signals-and-waves:ultrasound}{ultrageluid} van \qty{5.0}{MHz} in zacht weefsel
($v = \qty{1540}{m/s}$): wat is de \omterm{def:g12:mechanical-waves:wavelength}{golflengte}? \omterm{rem:g10:signals-and-waves:honest}{Echo}'s komen na
\qty{26}{\micro s} en \qty{58}{\micro s} terug van de twee wanden van een
orgaan: bepaal hun diepten en de dikte van het orgaan.
\end{exercise}

\begin{exercise}[$\star\star\star$]\label{exo:g12:sound-acoustics:13}
Een openluchtpodium levert \qty{110}{dB} op \qty{3.0}{m}. Op welke
afstand zakt het niveau tot \qty{85}{dB} als je het podium als kleine bron
behandelt? Geef twee redenen waarom de werkelijke afstand anders is.
\end{exercise}

\begin{exercise}[$\star\star\star$]\label{exo:g12:sound-acoustics:14}
Bepaal de verhouding van de pijngrens (\qty{1}{W/m^2}) tot $I_0$. Neem
het trommelvlies als \qty{0.5}{cm^2} en bereken het ontvangen \omterm{def:g11:work-of-force:power}{vermogen} bij
elke grens; becommentarieer het oor als detector.
\end{exercise}

\begin{exercise}[$\star\star\star$]\label{exo:g12:sound-acoustics:15}
Elke identieke festivalluidspreker geeft in zijn eentje \qty{80}{dB} bij
de mengtafel. Welk niveau geven er $2$? En $10$? En $100$? Wat kost elke
extra \qty{+10}{dB} --- en wat zegt dat over wedstrijdjes in \omterm{def:g12:sound-acoustics:pitch-loudness}{luidheid}?
\end{exercise}

\section{Probleem: Het atelier van de vioolbouwer}

\begin{problem}[{Weekendopgave --- het atelier van de vioolbouwer: een
gitaarsnaar fret voor fret uitgezet, haar boventonen tot een klankkleur
gestemd, een recital onder de gevarengrens gehouden, en een jacht met
ultrageluid op een fout in het hout}]
\label{pb:g12:sound-acoustics:1}
Een vioolbouwer maakt een gitaar af. De a-snaar heeft een trillende lengte
$L = \qty{65.0}{cm}$ en moet op $f_1 = \qty{110}{Hz}$ klinken. Neem
\qty{340}{m/s} voor \omterm{def:g10:signals-and-waves:sound}{geluid} in lucht, en neem $f_1 = v/(2L)$ aan
(\cref{prop:g12:sound-acoustics:string-modes}), met $v$ de golfsnelheid op
de snaar.

\medskip
\noindent\textbf{Deel I --- De toets uitzetten.}
\begin{enumerate}
  \item Bereken de golfsnelheid $v$ op de snaar.
  \item Bereken de \omterm{def:g12:mechanical-waves:wavelength}{golflengte} van de \omterm{def:g12:sound-acoustics:harmonics}{grondtoon} \emph{op de snaar}, en
        daarna de \omterm{def:g12:mechanical-waves:wavelength}{golflengte} van het \omterm{def:g10:signals-and-waves:sound}{geluid} van \qty{110}{Hz} \emph{in de
        lucht}. Waarom verschillen die bij dezelfde \omterm{def:g10:signals-and-waves:frequency}{frequentie}?
  \item Het octaaf (\qty{220}{Hz}): welke trillende lengte levert dat op?
        Waar moet de twaalfde fret zitten, als deel van $L$?
  \item De kwint (\qty{165}{Hz}): welke trillende lengte levert die op?
  \item Elke fret tilt de noot een halve toon op, een factor $2^{1/12}$.
        Ga met logaritmen na hoeveel halve tonen er tussen \qty{110}{Hz}
        en \qty{220}{Hz} liggen; welke lengte hoort bij de eerste fret?
\end{enumerate}

\medskip
\noindent\textbf{Deel II --- \omterm{def:g12:sound-acoustics:harmonics}{Boventonen} en \omterm{def:g12:sound-acoustics:timbre}{klankkleur}.}
\begin{enumerate}[resume]
  \item Noem de eerste vijf \omterm{def:g12:sound-acoustics:harmonics}{boventonen} van de losse a-snaar.
  \item De a van een violist klinkt op \qty{440}{Hz}. Welke \omterm{def:g12:sound-acoustics:harmonics}{boventoon} van
        de gitaarsnaar is dat, en welk muzikaal interval scheidt hem van
        de \omterm{def:g12:sound-acoustics:harmonics}{grondtoon}? En \qty{330}{Hz}?
  \item De bouwer tokkelt eerst bij de kam, daarna boven de toets:
        dezelfde \omterm{def:g12:sound-acoustics:pitch-loudness}{toonhoogte}, een andere kleur. Verklaar dat met spectra.
  \item Een lichte aanraking in het midden tijdens het tokkelen dwingt
        daar een \omterm{def:g12:sound-acoustics:standing-wave}{knoop} af. Welke \omterm{def:g12:sound-acoustics:harmonics}{boventonen} overleven, en welke
        \omterm{def:g12:sound-acoustics:pitch-loudness}{toonhoogte} hoor je?
  \item Hoeveel \omterm{def:g12:sound-acoustics:harmonics}{boventonen} van de losse snaar liggen in beginsel binnen
        het hoorbare gebied?
\end{enumerate}

\medskip
\noindent\textbf{Deel III --- Het recital.}
Bij het eerste concert van de gitaar straalt elk instrument een
akoestisch \omterm{def:g11:work-of-force:power}{vermogen} $P = \qty{6.0e-3}{W}$ uit als kleine bron, in alle
richtingen even sterk.
\begin{enumerate}[resume]
  \item Bereken de intensiteit op $r = \qty{3.0}{m}$ van één instrument.
  \item Leid daaruit het \omterm{def:g12:sound-acoustics:level}{geluidsniveau} daar af.
  \item Er speelt een kwartet: welk niveau geven vier zulke instrumenten
        op dezelfde plek? Waarom is het niet vier keer het niveau?
  \item Er spelen twaalf instrumenten: welk niveau geldt op \qty{3.0}{m}?
  \item Hoever moet een toehoorder achteruit gaan opdat het ensemble tot
        \qty{85}{dB} zakt?
  \item Veiligheidsregel: \qty{85}{dB} is acht uur veilig, en elke
        \qty{+3}{dB} halveert de veilige tijd. Hoe lang mag de toehoorder
        van vraag 14 op de eerste rij veilig blijven zitten?
\end{enumerate}

\medskip
\noindent\textbf{Deel IV --- De fout in het hout.}
Vóór het vernissen tast de bouwer een halsblok van \qty{45}{mm} dik af met
een tester van \qty{2.0}{MHz}; neem in dit hout $v = \qty{5.0e3}{m/s}$.
\begin{enumerate}[resume]
  \item Bereken de \omterm{def:g12:mechanical-waves:wavelength}{golflengte} van het \omterm{def:g10:signals-and-waves:ultrasound}{ultrageluid} in het hout. Waarom
        moet een \omterm{def:g10:signals-and-waves:wave}{golf} die fouten opspoort ultrasoon zijn en niet hoorbaar?
  \item Bereken de echovertraging van het achtervlak van het blok.
  \item Er komt een \omterm{rem:g10:signals-and-waves:honest}{echo} terug na \qty{7.2}{\micro s}: hoe diep zit de
        fout?
  \item Schrijf de werkplaatskaart uit: golfsnelheid op de snaar, lengten
        bij de twaalfde en de eerste fret, veilige luisterafstand en
        diepte van de fout.
\end{enumerate}
\end{problem}
